\bookchapter{Basic function shapes used in AI: activation functions}\label{ch:activation}

\section*{Explanation}
A neuron first computes a weighted sum $z$ and then applies an \emph{activation function} $a(z)$. Why is this needed? If all functions were linear, the composition of two of them would be linear again:
\[
f(g(x))=a_1(a_2x+b_2)+b_1=(a_1a_2)\,x+(a_1b_2+b_1).
\]
A network made only of linear functions could therefore never be more than one single line. Activation functions \emph{bend} the graph and give networks their power.

\begin{center}
\begin{tikzpicture}
\begin{groupplot}[group style={group size=2 by 2, horizontal sep=2cm, vertical sep=2cm},
  width=6.8cm,height=4.4cm,axis lines=middle,xlabel={$z$},
  grid=both, grid style={line width=0.1pt, draw=gray!30}, major grid style={line width=0.2pt, draw=gray!50},
  tick label style={font=\footnotesize}, title style={font=\small\bfseries},
  every axis plot/.append style={thick,blue}, samples=100]
  \nextgroupplot[title={Step}, xmin=-4,xmax=4,ymin=-0.5,ymax=1.5, xtick={-4,-2,...,4}, ytick={0,1}]
    \addplot[domain=-4:0,samples=2]{0};
    \addplot[domain=0:4,samples=2]{1};
    \addplot[only marks,mark=*,blue] coordinates {(0,1)};
    \addplot[only marks,mark=o,blue] coordinates {(0,0)};
  \nextgroupplot[title={Sigmoid $\sigma(z)=\frac{1}{1+e^{-z}}$}, xmin=-6,xmax=6,ymin=-0.3,ymax=1.3, xtick={-6,-3,...,6}, ytick={0,0.5,1}]
    \addplot[domain=-6:6]{1/(1+exp(-x))};
  \nextgroupplot[title={ReLU $\max(0,z)$}, xmin=-4,xmax=4,ymin=-1,ymax=4, xtick={-4,-2,...,4}, ytick={0,1,...,4}]
    \addplot[domain=-4:4,samples=3]{max(0,x)};
  \nextgroupplot[title={tanh $z$}, xmin=-4,xmax=4,ymin=-1.5,ymax=1.5, xtick={-4,-2,...,4}, ytick={-1,0,1}]
    \addplot[domain=-4:4]{tanh(x)};
\end{groupplot}
\end{tikzpicture}

\captionof{figure}{The four classic activation functions.}\label{fig:b1_10}
\end{center}

\begin{center}
\renewcommand{\arraystretch}{1.7}
\begin{tabular}{llll}
\textbf{Name} & \textbf{Formula} & \textbf{Range} & \textbf{Remark}\\\hline
Step & $0$ for $z<0$, $1$ for $z\ge0$ & $\{0,1\}$ & first artificial neuron, not smooth\\
Sigmoid & $\dfrac1{1+e^{-z}}$ & $(0,1)$ & smooth, output like a probability\\
ReLU & $\max(0,z)$ & $[0,\infty)$ & very cheap, most popular today\\
tanh & $\dfrac{e^z-e^{-z}}{e^z+e^{-z}}$ & $(-1,1)$ & like sigmoid, but centred at $0$
\end{tabular}
\end{center}

Useful values: $\sigma(0)=0.5$, $\sigma(2)\approx0.881$, $\sigma(-2)\approx0.119$ (in general $\sigma(-z)=1-\sigma(z)$), $\tanh(1)\approx0.762$.

\begin{example}
A neuron with inputs $x=(2,1)$, weights $w=(1,-2)$ and bias $b=0.5$ has $z=2-2+0.5=0.5$. Its output is
\[
\text{step}: 1,\qquad \sigma(0.5)\approx0.622,\qquad \text{ReLU}: 0.5,\qquad \tanh(0.5)\approx0.462.
\]
\end{example}

\begin{aibox}
The slopes of these curves are needed in training: in Chapter~\ref{ch:keyfunctions} you will see that the derivative of the sigmoid is $\sigma'(z)=\sigma(z)(1-\sigma(z))$ and that ReLU has slope $0$ for $z<0$ and $1$ for $z>0$. Notice how the sigmoid is almost flat for large $|z|$: there the slope is nearly $0$ and learning is slow.
\end{aibox}

\section*{Practice problems}
\begin{enumerate}[label=\textbf{\thechapter.\arabic*}]
  \item Compute step$(-3)$, step$(0)$, step$(2)$, ReLU$(-2)$ and ReLU$(3.5)$.
  \item Compute $\sigma(0)$, $\sigma(2)$ and $\sigma(-2)$. What do you notice about $\sigma(2)+\sigma(-2)$?
  \item Which of the four activation functions can output a negative number? Which can output a number larger than $1$?
  \item A neuron has inputs $x=(1,2,-1)$, weights $w=(0.5,-0.25,0.2)$ and bias $b=0.1$. Compute $z$ and the output for the step, sigmoid, ReLU and tanh activation.
  \item Show that the composition of two linear functions $f(x)=a_1x+b_1$ and $g(x)=a_2x+b_2$ is again linear. What are its slope and intercept?
  \item Looking at Figure~\ref{fig:b1_10}, for which $z$ is the sigmoid almost flat? What does that mean for the neuron's output when $z=10$?
  \item $\star$ Show that $\tanh(z)=2\sigma(2z)-1$.
  \item \texttt{[code]} Plot all four activation functions with Python (matplotlib) for $z\in[-5,5]$.
\end{enumerate}

\newpage
